% TOP_PROBLEM: {"id": 3700001, "title": "Fuchsian equations with unitary monodromy", "category_id": 8, "category": {"id": 8, "name": "analysis", "display_name": "Analysis", "description": "Limits, continuity, calculus, and function theory.", "slug": "analysis", "order_index": 8, "created_at": "2026-07-31T15:26:25.671Z"}, "status": "partially_solved", "problem_number": "AMR-036-0001", "set_id": 13, "statusQualification": "An explicit projective normal form defines the accessory parameters; the source's established four-singularity case is separated."}
% FIRST_POSTED: 2026-10-01
% ENTRY_KIND: statement_only
% UnsolvedMath ID 3700001; source catalogue code AMR-036-0001
% !TeX program = lualatex
% Definition and sources only; this file is not an LLM solution attempt.
\documentclass[11pt,a4paper]{article}
\usepackage[margin=25mm]{geometry}
\usepackage{fontspec}
\setmainfont{lmroman10-regular.otf}[BoldFont=lmroman10-bold.otf,
 ItalicFont=lmroman10-italic.otf,BoldItalicFont=lmroman10-bolditalic.otf]
\usepackage{amsmath,amssymb,mathtools}
\usepackage{unicode-math}
\setmathfont{latinmodern-math.otf}
\AtBeginDocument{\let\setminus\smallsetminus}
\usepackage{xurl,hyperref}
\hypersetup{colorlinks=true,urlcolor=blue}
\setcounter{secnumdepth}{0}
\setlength{\parindent}{0pt}
\setlength{\parskip}{5pt}
\setlength{\emergencystretch}{3em}
\begin{document}
\section*{Fuchsian equations with unitary monodromy}\label{3700001}
{\small UnsolvedMath ID 3700001 · AMR-036-0001 · Analysis · AMR Open Problem Lists}
\subsection{Definitions and mathematical statement}
Prescribe $n\geq4$ distinct points on the Riemann sphere and real
exponent differences $\alpha_1,\ldots,\alpha_n$. After a change of
coordinate put the last point at infinity and write the others
as $a_1,\ldots,a_{n-1}\in\mathbb C$. In projective normal form the
second-order equation is
\[
 y''+Q(z)y=0,\qquad
 Q(z)=\sum_{j=1}^{n-1}
       \left(\frac{1-\alpha_j^2}{4(z-a_j)^2}
                         +\frac{\beta_j}{z-a_j}\right).
\]
The accessory parameters $\beta_j$ satisfy
\[
 \sum_j\beta_j=0,\qquad
 \sum_j\left(\frac{1-\alpha_j^2}{4}+a_j\beta_j\right)
                         =\frac{1-\alpha_n^2}{4}.
\]
These conditions impose the prescribed regular-singular behavior
at infinity, leaving an affine parameter space of dimension $n-3$.
The exponent difference is the difference of the two local
indicial roots; its sign does not affect this normal form.

Analytic continuation of the ratio of two independent local
solutions around loops avoiding the marked points acts by
Möbius transformations. This is the projective monodromy group,
defined up to conjugacy. Let $E$ be the accessory parameters
for which this group is conjugate into
$\operatorname{PSU}(2)=\operatorname{SU}(2)/\{\pm I\}$, the
rotation group of the round sphere.

For every prescribed configuration and real exponent data, is
$E$ discrete in the accessory parameter space? Is it finite?
Discreteness means every point of $E$ is isolated; finiteness
also excludes infinitely many isolated parameters escaping to
infinity. The source establishes finiteness for $n=4$; the
general remaining question concerns $n\geq5$.
\subsection{Short English statement}
For fixed singularities and exponent differences, are there only finitely many Fuchsian equations with spherical unitary monodromy?
\subsection{Sources}
\begin{itemize}
\item[S1] ULAM.AI and the UnsolvedMath Contributors, supplied problems.json, record 3700001 (AMR-036-0001), AMR Open Problem Lists. Catalogue identity and source transcription; CC BY 4.0.
\url{https://huggingface.co/datasets/ulamai/UnsolvedMath}.
\item[S2] Eremenko - My favorite unsolved problems, Fuchsian equations, two questions. Original source indexed by AMR-036-0001.
\url{https://www.math.purdue.edu/~eremenko/uns1.html}.
\item[S3] A. Eremenko, Fuchsian equations with unitary monodromy. Author-maintained primary problem note.
\url{https://www.math.purdue.edu/~eremenko/dvi/fuchsian.pdf}.
\end{itemize}
\end{document}
